\documentclass[a4paper,12pt]{article}

% Math packages
\usepackage{amsmath, amssymb, amsthm}

% Header
\usepackage{fancyhdr}
\pagestyle{fancy}

\fancyhead{} % clear default
\fancyhead[L]{Anton Thestrup Jørgensen}
\fancyhead[R]{s194268}

\begin{document}

\section*{Question 1}

\subsection*{a)}

The set is closed under the group operation:
\begin{align*}
    f_{a, b} \circ f_{c, d}[x] &= f_{a, b}[f_{c, d}[x]] \\
    &= f_{a, b}[cx + d] \\
    &= a(cx + d) + b \\
    &= acx + ad + b \\
    &= (ac)x + (ad + b) \\
    &= f_{ac, ad + b}[x]
\end{align*}

The operation is associative:
\begin{align*}
    (f_{a,b} \circ f_{c,d}) \circ f_{e, g}[x] &= f_{a,b} \circ f_{c,d}[ex + g] \\
    &= f_{ac, ad + b}[ex + g] \\
    &= ac(ex + g) + ad + b \\
    &= acex + acg + ad + b \\
    &= a(cex + cg + d) + b \\
    &= f_{a,b}[cex + cg + d] \\
    &= f_{a,b} \circ f_{ce, cg + d}[x] \\
    &= f_{a, b} \circ (f_{c, d} \circ f_{e, g})[x]
\end{align*}

There is an identity:
\begin{align*}
    f_{1, 0} \circ f_{a,b} &= 1(ax + b) + 0 = ax + b = f_{a, b} \\
    f_{a, b} \circ f_{1, 0} &= a(1x + 0) + b = ax + b = f_{a, b}
\end{align*}

All elements have an inverse:
\begin{align*}
    f_{a^{-1}, -a^{-1}b} \circ f_{a,b} &= a^{-1}(ax + b) - a^{-1}b = x + a^{-1}b -a^{-1}b = x = f_{1,0} \\
    f_{a, b} \circ f_{a^{-1}, -a^{-1}b} &= a(a^{-1}x - a^{-1}b) + b = x - b + b = x = f_{1, 0}
\end{align*}\\

Therefore $(G, \circ)$ is a group.

\subsection*{b)}

$T$ is a subgroup, since it satisfies the three criterions of the definitions of a subgroup:

\begin{enumerate}
    \item $e = f_{1, 0} \in T$
    \item $f_{1, b}^{-1} = f_{1, -b} \in T$
    \item $f_{1, b} \circ f_{1, d} = f_{1, d} \circ f_{1, b} = f_{1, b + d} \in T$
\end{enumerate}
In the inclusion statements it is used that $0$, $-b$, and $b + d$ respectively are real numbers. It is also taken as `obvious' that $T$ is a subset of $G$.

\subsection*{c)}
For a fixed $g = f_{a, b}$, the elements of $gT$ are on the form:
$$
f_{a, b} \circ f_{1, d} = a(x + d) + b = ax + ad + b = ax + (ad + b)
$$
Likewise, elements of $Tg$ are on the form:
$$
f_{1, d} \circ f_{a, b} = 1(ax + b) + d = ax + b + d = ax + (b + d)
$$

As $d$ varies over all real numbers when generating the left and right cosets, so does $ad + b$ and $b + d$ respectively. Therefore both $gT$ and $Tg$ are equal to the set of lines with slope $a$, and by transitivity to each other.

The equality could also be proven using the standard double inclusion tactic.

\clearpage
\section*{Question 2}


\subsection*{a)}

If we can find an element of order $2p$, we have shown that $G$ is cyclic. By Lagrange's Theorem, we have $\forall g \in G, \; ord(g) \; \text{divides} \; G$. This means that $ord(g) \in \{1, 2, p, 2p\}$. Now choose an element $g \in G \setminus \{e, \alpha \}$ (which is not empty, since the order of $G$ is $2p \geq 6$). Since the only element of order $1$ is $e$, and the only element of order $2$ is $\alpha$, $g$ must have order $p$ or $2p$. If it has order $2p$, $G$ is cyclic.

Consider then that $g$ has order $p$, look at the element $g \alpha g^{-1}$, and see that $(g \alpha g^{-1})^2 = e$. This means that $ord(g \alpha g^{-1}) = 2$ (normally we would show that $2$ is the smallest such number, but for $2$ it is trivial). But $\alpha$ is the only element of order $2$, so $g \alpha g^{-1} = \alpha$. But this means that $g \alpha = \alpha g$, so $g$ and $\alpha$ commute!

Now take $\alpha g$ raised to successive powers:

\begin{align*}
    (\alpha g)^1 &= \alpha g \\
    (\alpha g)^2 &= g^2 \\
    (\alpha g)^3 &= \alpha g^3 \\
    \dots \\
    (\alpha g)^p &= \alpha \\
    (\alpha g)^{p + 1} &= g \\
    (\alpha g)^{p + 2} &= \alpha g^2 \\
    \dots \\
    (\alpha g)^{2p} &= e
\end{align*}

There are no identity elements 'hiding' in the $\dots$ sections, since $\alpha$ is never the inverse of $g^k$ (inverses are unique, $\alpha$ is its own inverse, and by order of powers $ord(g^k) \neq 2$). Then $2p$ is the smallest $k$ such that $(\alpha g)^k = e$, and $ord(\alpha g) = 2p$. Therefore $G$ is cyclic!

\newpage
\subsection*{b)}

The actions of $0$ and $1$ are given by definition and in the problem statement. Otherwise the strategy will be to derive the action using permutation, and when the action of an inverse is known, to use the fact that $\varphi(g^{-1})^{-1} = \varphi(g)$:

\begin{align*}
\varphi(0) &:= \mathrm{id}_A \\
\varphi(1) &:= (3\ 4\ 5) \\
\varphi(2) &= \varphi(1+_6 1) = \varphi(1) \circ \varphi(1) = (3\ 5\ 4) \\
\varphi(3) &= \varphi(3^{-1})^{-1} = \varphi(3)^{-1} \rightarrow \varphi(3) = \mathrm{id}_A \\
\varphi(4) &= \varphi(4^{-1})^{-1} = \varphi(2)^{-1} = (4\ 5\ 3) = (3\ 4\ 5) \\
\varphi(5) &= \varphi(5^{-1})^{-1} = \varphi(1)^{-1} = (5\ 4\ 3) = (3\ 5\ 4)
\end{align*}

It could also be used that since $\varphi(3) = \mathrm{id}_A$, $\varphi(4) = \varphi(3 + 1) = \mathrm{id}_A \circ \varphi(1) = \varphi(1)$ and so on. 

\end{document}